\clearpage
\section{Problem 5: Fan-out Event Distribution atau Publish--Subscribe}
\label{sec:solusi_p5}

\subsection{Rancangan Solusi dan Mekanisme}
\label{subsec:p5_rancangan}
Pada sistem agregasi bencana multi-lembaga, ragam pemangku kepentingan hilir terus bertambah seiring waktu (pembaruan metrik visual pada \textit{dashboard}, pengiriman notifikasi darurat ke tim lapangan, serta portal pemda eksternal).

Apabila \textit{Aggregator Service} memanggil seluruh sistem hilir secara sinkron melalui HTTP/REST langsung (\textit{point-to-point synchronous chaining}), muncul permasalahan \textit{cascading latency} dan \textit{tight coupling}: keterlambatan pada satu konsumen akan menahan siklus eksekusi \textit{Aggregator} dan menunda pengiriman data ke konsumen lain yang mendesak, serta penambahan konsumen baru mewajibkan modifikasi kode dan \textit{redeploy}.

Untuk mengatasi permasalahan tersebut, perancangan sistem menerapkan pola arsitektur \textbf{\textit{Publish--Subscribe}} berbasis \textit{log-based message broker} menggunakan Redpanda (kompatibel penuh dengan protokol Apache Kafka). Rancangan dan mekanisme solusi adalah sebagai berikut:
\begin{enumerate}
    \item \textbf{Decoupled Event Publishing:}
    Setelah \textit{Aggregator Service} menerima, memetakan, dan memvalidasi rekaman bencana dari lembaga sumber, \textit{Aggregator} menerbitkan entitas kanonikal \texttt{HazardEvent} ke \textit{topic} \texttt{bnpb.hazard-events} pada Redpanda. \textit{Topic} ini dikonfigurasi dengan 3 partisi dan memanfaatkan \texttt{hazard\_id} sebagai \textit{partition key} untuk menjamin \textit{causal ordering} per kejadian bencana. Penerbitan dilakukan secara asinkron tanpa keterikatan terhadap identitas, alamat IP, maupun ketersediaan konsumen mana pun.
    \item \textbf{Independent Consumer Groups:}
    Setiap layanan konsumen hilir beroperasi di bawah \textit{consumer group} yang terisolasi secara independen:
    \begin{itemize}
        \item \texttt{bnpb-dashboard-group}: Dikelola oleh \textit{Dashboard Consumer} untuk agregasi metrik frekuensi dan keparahan bencana.
        \item \texttt{bnpb-notifier-group}: Dikelola oleh \textit{Notifier Consumer} untuk mengirimkan notifikasi simulasi peringatan tingkat \textit{Waspada}, \textit{Siaga}, dan \textit{Awas}.
        \item \texttt{bnpb-evaluator-audit-group} / Konsumen Eksternal: Konsumen ketiga (misalnya audit atau portal pemda) yang dapat mendaftar dan membaca aliran data secara independen.
    \end{itemize}
    Masing-masing kelompok konsumen mengelola posisi penunjuk pembacaan (\textit{committed offset}) secara terpisah pada metadata Redpanda (\texttt{\_\_consumer\_offsets}). Kecepatan konsumsi maupun gangguan pada satu grup tidak berdampak sama sekali terhadap grup lainnya.
    \pagebreak
    \item \textbf{Consumer Downtime Resilience dan Backlog Recovery:}
    Apabila suatu layanan konsumen (misalnya \textit{Dashboard}) mengalami kegagalan proses atau dihentikan (\textit{downtime}), Redpanda tetap menyimpan pesan-pesan baru yang diterbitkan produser ke dalam \textit{append-only log} pada disk broker. Nilai \textit{consumer lag} untuk grup tersebut akan bertambah secara proporsional. Sementara itu, konsumen lain (\textit{Notifier}) tetap memproses pesan baru secara \textit{real-time} tanpa terganggu. Ketika layanan yang mati dihidupkan kembali, konsumen membaca ulang konfigurasi \textit{committed offset} terakhirnya dan segera mengejar ketertinggalan (\textit{catch-up backlog}) hingga nilai \textit{lag} kembali ke nol.
    \item \textbf{Manual Offset Commit Discipline:}
    Konsumen mengimplementasikan pustaka \texttt{kafka-go} dengan pembacaan berbasis \texttt{FetchMessage}. \textit{Offset} pesan tidak dikomit secara otomatis, melainkan secara manual melalui \texttt{CommitMessages} tepat setelah pemrosesan data internal dan deduplikasi lokal berhasil diselesaikan. Hal ini menjamin semantik pengiriman \textit{at-least-once}.
\end{enumerate}

\subsection{Alasan Keputusan dan Trade-Off}
\label{subsec:p5_alasan}
Pemilihan pola dan teknologi distribusi \textit{event} didasarkan pada perbandingan arsitektur berikut:

\begin{enumerate}
    \item \textbf{Pola Publish--Subscribe vs. Pola Polling Sinkron / Webhook:}
    \begin{itemize}
        \item \textit{Polling Terjadwal:} Jika konsumen secara berkala melakukan \mbox{HTTP} GET ke \textit{Aggregator}, timbul beban kueri basis data berulang (\textit{database thrashing}) dan latensi deteksi bencana menjadi tinggi tergantung interval \textit{polling}.
        \item \textit{Direct Webhook/HTTP Calls:} Membebani \textit{Aggregator} dengan tanggung jawab \textit{retry}, manajemen kegagalan koneksi, \textit{timeout handling}, dan registrasi dinamis URL konsumen.
        \item \textit{Keputusan:} Pola \textit{Publish--Subscribe} dipilih karena memisahkan tanggung jawab secara bersih (\textit{separation of concerns}). Produser berfokus pada validasi dan penerbitan \textit{event}, sedangkan broker menangani retensi, \textit{buffering}, partisi, dan distribusi data.
    \end{itemize}
    \item \textbf{Pemilihan Redpanda (Log-based Pub/Sub) vs. Tradisional Message Queue (RabbitMQ AMQP):}
    \begin{itemize}
        \item \textit{RabbitMQ (AMQP Fanout Exchange):} Menyediakan mekanisme distribusi pesan yang baik, namun model antreannya menghapus pesan setelah diakui (\textit{destructive read}). Apabila konsumen baru mendaftar di kemudian hari atau konsumen yang mati membutuhkan pemutaran ulang (\textit{replayability}) dari \textit{offset} awal hari, antrean AMQP membutuhkan konfigurasi antrean terpisah sebelum pesan diterbitkan.
        \item \textit{Redpanda (Kafka API):} Menggunakan model \textit{immutable distributed commit log}. Pesan dipertahankan berdasarkan durasi retensi waktu dan kuota penyimpanan fisik, bukan status pembacaan konsumen. Konsumen baru dapat bergabung kapan saja dan membaca data tanpa intervensi produser. Redpanda dipilih karena diimplementasikan dalam C++ tanpa ketergantungan JVM atau Zookeeper, dan kompatibel dengan API Kafka. Hemat memori bukan alasan pemilihan: saat idle, Redpanda (352 MiB) tidak lebih ringan daripada Kafka 3.7 mode KRaft (297 MiB; lihat ADR 002).
    \end{itemize}
\end{enumerate}

\clearpage
\noindent\textbf{Trade-Off Arsitektural:}
\begin{itemize}
    \item \textbf{Kelebihan:} Produser tidak bergantung pada identitas atau ketersediaan konsumen; broker mengisolasi gangguan komponen hilir (\textit{fault isolation}); dan konsumen baru dapat ditambahkan tanpa mengubah produser.
    \item \textbf{Konsekuensi Negatif:} Diperlukan pengelolaan infrastruktur broker pesan tambahan (\textit{cluster health}, partisi, dan kuota disk); data pada sistem hilir bersifat \textit{eventually consistent} dan dapat tertunda ketika terdapat backlog; konsumen wajib menangani potensi duplikasi pesan akibat semantik \textit{at-least-once}. Latensi \textit{fan-out} belum diukur secara khusus pada Milestone 1.
\end{itemize}

\subsection{Batasan yang Disadari}
\label{subsec:p5_batasan}
Batasan teknis yang disadari dan sengaja dibatasi pada lingkup implementasi saat ini meliputi:
\begin{enumerate}
    \item \textbf{Deduplikasi Dalam Memori Bersifat Ephemeral:}
    Pada implementasi prototipe \textit{Dashboard} dan \textit{Notifier}, deduplikasi \texttt{hazard\_id} disimpan dalam struktur \textit{hash map} berbasis memori (\texttt{sync.Mutex} di Go). Apabila kontainer konsumen dihentikan secara paksa atau di-\textit{restart}, riwayat kunci deduplikasi dalam memori akan tereset.
    \item \textbf{Edge Case At-Least-Once pada Crash Antara Eksekusi dan Commit:}
    Apabila konsumen mengalami kegagalan sistem (\textit{panic} atau \textit{SIGKILL}) tepat setelah efek samping dijalankan (misalnya mencetak log agregasi) tetapi sebelum fungsi \texttt{reader.CommitMessages()} dieksekusi, pesan yang sama akan dikonsumsi kembali setelah kontainer hidup kembali.
    \item \textbf{Batas Retensi Log Broker:}
    Penyimpanan pesan pada Redpanda dibatasi oleh kebijakan retensi waktu dan ukuran log. Jika konsumen mengalami \textit{downtime} melebihi durasi retensi log yang dikonfigurasi, pesan-pesan terlama yang belum terbaca akan terhapus secara otomatis oleh broker sehingga konsumen tidak dapat lagi memulihkan data tersebut secara utuh.
\end{enumerate}

\clearpage
\subsection{Bukti Penyelesaian}
\label{subsec:p5_bukti}
Pengujian dan validasi solusi Problem 5 dilakukan secara otomatis melalui skrip \path{scripts/test_consumer_recovery.sh} (dapat dijalankan melalui \texttt{make demo-p5}). Seluruh pengujian mencakup verifikasi partisi, penerbitan paralel, pemulihan \textit{downtime}, dan penambahan konsumen dinamis.

\noindent\begin{minipage}{\linewidth}
\begin{lstlisting}[language=bash, caption={Keluaran Eksekusi Pengujian Otomatis Problem 5 (Fan-Out \& Downtime Recovery)}]
[Step 0] Checking Redpanda topic: 'bnpb.hazard-events' (3 partitions) -> [PASS]
[Step 1] Starting Dashboard and Notifier containers -> [PASS]
[Step 2] Publishing event 'c08-fanout-...' -> Fan-out received by both -> [PASS]
[Step 3] Stopping Dashboard and publishing recovery event during downtime:
GROUP                 TOPIC               STATUS    TOTAL-LAG
bnpb-dashboard-group  bnpb.hazard-events  Empty     1
[PASS] Notifier processed; Dashboard offline and accumulated TOTAL-LAG = 1
[Step 4] Restarting Dashboard and verifying backlog recovery:
GROUP                 TOPIC               STATUS    TOTAL-LAG
bnpb-dashboard-group  bnpb.hazard-events  Stable    0
[PASS] Dashboard resumed committed offset; TOTAL-LAG returned to zero
[Step 5] Dynamic 3rd consumer joined -> Events received -> [PASS]
Results: 11 passed, 0 failed. Fan-Out & Recovery Verified!
\end{lstlisting}
\end{minipage}

\noindent\textbf{Analisis Bukti Log:}
\begin{enumerate}
    \item \textbf{Fan-Out Verifikasi (Langkah 2):} Satu kejadian bencana (\texttt{c08-fanout-...}) diterbitkan ke Redpanda. Kedua layanan (\textit{Dashboard} dan \textit{Notifier}) mendeteksi dan memproses pesan tersebut secara simultan dalam kelompok konsumen masing-masing (\texttt{bnpb-}\allowbreak\texttt{dashboard-}\allowbreak\texttt{group} dan \texttt{bnpb-}\allowbreak\texttt{notifier-}\allowbreak\texttt{group}).
    \item \textbf{Downtime \& Lag Accumulation (Langkah 3):} Saat kontainer \textit{Dashboard} dimatikan paksa, pengiriman event baru (\texttt{c08-recovery-...}) tetap diproses normal oleh \textit{Notifier}. Pemeriksaan Redpanda via \texttt{rpk group describe} menunjukkan \texttt{TOTAL-LAG = 1} pada \textit{Dashboard}, membuktikan pesan aman tersimpan pada broker dan tidak hilang.
    \item \textbf{Backlog Catch-Up (Langkah 4):} Begitu kontainer \textit{Dashboard} dinyalakan kembali, konsumen segera menyerap pesan yang tertinggal dan melakukan \textit{manual commit}, ditandai dengan \texttt{TOTAL-LAG} yang kembali menjadi 0.
    \item \textbf{Penambahan Konsumen Baru (Langkah 5):} Kontainer independen ketiga (\texttt{bnpb-}\allowbreak\texttt{evaluator-}\allowbreak\texttt{audit-...}) dijalankan pada topik \texttt{bnpb.hazard-events} tanpa melakukan \textit{restart} atau pengubahan kode apa pun pada \textit{Aggregator Service}. Konsumen baru langsung menerima \textit{event} secara transparan.
\end{enumerate}

\clearpage
\subsection{Jawaban Butir ``Di Laporan''}
\label{subsec:p5_jawaban}
Jawaban komprehensif atas aspek evaluasi Problem 5:

\begin{enumerate}
    \item \textbf{Broker dan Pola yang Dipilih:}
    \begin{itemize}
        \item \textbf{Broker (Redpanda):} Kompatibel API Kafka v3, arsitektur \textit{single binary} C++ (tanpa JVM dan ZooKeeper), serta dukungan utilitas CLI \texttt{rpk}.
        \item \textbf{Pola (Pub--Sub Log Terdistribusi):} Topik \texttt{bnpb.hazard-events} dibagi ke 3 partisi dengan \textit{consumer groups} mandiri untuk replikasi aliran pesan (\textit{fan-out}) ke setiap grup secara otonom.
    \end{itemize}

    \item \textbf{Semantik Pengiriman yang Berlaku (\textit{At-Least-Once Delivery}):}
    \begin{itemize}
        \item Produser (outbox relay pada Aggregator) mengirimkan pesan dengan tingkat pengakuan \texttt{acks=1} (\texttt{kafka.RequireOne}), sehingga pesan telah dicatat oleh \textit{leader} partisi sebelum dianggap sukses.
        \item Konsumen tidak memakai komit otomatis: pesan dibaca via \texttt{reader.FetchMessage}, lalu memanggil \texttt{reader.CommitMessages} secara manual tepat setelah pemrosesan dan deduplikasi selesai.
        \item Jika terjadi kegagalan jaringan atau proses mati sebelum komit, \textit{offset} belum tersimpan sehingga pesan dapat dikirim ulang (\textit{redeliver}) saat pulih. Kombinasi transactional outbox, retry, retensi broker, dan komit manual mendukung semantik \textit{at-least-once} selama basis data dan log broker tetap tersedia; rancangan single-node dengan \texttt{acks=1} bukan jaminan absolut \textit{zero message loss}.
    \end{itemize}

    \item \textbf{Konsekuensi bagi Consumer dan Perlunya Sifat Idempoten:}
    \begin{itemize}
        \item \textbf{Kebutuhan Idempotensi:} Dalam semantik \textit{at-least-once}, pengiriman ganda (\textit{duplicate delivery}) dapat terjadi saat \textit{timeout}, retry, atau partisi \textit{rebalance}. Tanpa idempotensi, metrik bencana dapat terhitung ganda (\textit{double count}) pada Dashboard dan memicu notifikasi berulang (\textit{alert fatigue}) pada Notifier.
        \item \textbf{Strategi Mitigasi (\textit{Event Identity}):} Setiap bencana membawa identitas global deterministik (\texttt{hazard\_id}). Pada prototipe saat ini, konsumen mengelola set \texttt{processed[hazard\_id]} dalam memori dengan perlindungan \texttt{sync.Mutex} untuk mengabaikan duplikasi sebelum komit. Pada skala multi-replika produksi, hal ini disikapi via operasi atomik (Redis \texttt{SETNX}) atau integritas basis data relasional (\texttt{ON CONFLICT DO NOTHING}).
    \end{itemize}
\end{enumerate}
